% Figure from MATH 591 notes §85 (Lecture 34; not from lecture): the curve alpha_p(t) in the fixed space Lambda^k T_p^*M
\begin{tikzpicture}[>=Stealth]
\draw[thick, fill=gray!6] plot[smooth cycle] coordinates {(-0.3,-0.9) (1.6,-1.2) (3.8,-0.8) (4.7,0.7) (3.6,1.75) (0.8,1.5) (-0.4,0.6)};
\node[font=\small] at (3.95,-1.0) {$M$};
\draw[thick, red!70!black, ->] plot[smooth] coordinates {(0.4,-0.3) (1.3,0.1) (2.4,0.42) (3.4,0.25)};
\node[red!70!black, font=\scriptsize] at (1.45,-0.45) {$t \mapsto \varphi_t(p)$};
\fill (0.4,-0.3) circle (1.5pt) node[left, font=\scriptsize] {$p$};
\fill (2.4,0.42) circle (1.5pt) node[below=2pt, font=\scriptsize] {$\varphi_t(p)$};
\foreach \d in {-0.14,0,0.14} { \draw[blue!70!black, thick] ({2.55+\d},{0.72+\d*0.3}) -- ({2.8+\d},{1.08+\d*0.3}); }
\node[blue!70!black, font=\scriptsize, anchor=west] at (3.0,1.15) {$\alpha_{\varphi_t(p)}$};
\draw[->, dashed, blue!70!black] (2.35,1.0) to[bend right=35] node[pos=0.45, above left=-1pt, font=\scriptsize] {$(d\varphi_t)_p^*$} (0.45,-0.12);
\begin{scope}[xshift=6.4cm]
\fill[gray!12] (-0.3,-1.0) -- (3.2,-1.0) -- (3.8,1.4) -- (0.3,1.4) -- cycle;
\node[font=\scriptsize, gray!60!black, anchor=west] at (2.4,1.65) {$\Lambda^k T_p^*M$ (fixed)};
\draw[thick, orange!85!black] plot[smooth] coordinates {(0.2,-0.75) (0.6,-0.45) (1.6,0.2) (2.7,0.32) (3.2,0.75)};
\fill (0.6,-0.45) circle (1.2pt) node[left=1pt, font=\scriptsize] {$\alpha_p(-t)$};
\fill (1.6,0.2) circle (1.5pt) node[below=3pt, font=\scriptsize] {$\alpha_p(0) = \alpha_p$};
\fill (2.7,0.32) circle (1.2pt) node[below right=-1pt, font=\scriptsize] {$\alpha_p(t)$};
\draw[->, very thick, red!70!black] (1.6,0.2) -- (2.45,0.55) node[above, font=\scriptsize] {$(\mathcal{L}_\mathbf{X}\alpha)_p$};
\end{scope}
\end{tikzpicture}
